% c03 —— §4 末：格极化 K3 曲面与周期区域（p006–p007 / 原书 p.132–133）
% 承 c02 末 (4.14) 的「……とおく．」（译作「这样定义。」）；本块从定义 4.1 起，
% 止于 p007 末「……$\mathcal{D}/\Gamma$ のトロイダル部分コンパクト化が得られる．」句完。

% ===== 原书 p.132 =====

\begin{definition}[Dolgachev {[Dol96]}]
K3 曲面 $Y$ 与格的原始嵌入 $j\colon M\hookrightarrow\Pic(Y)$ 组成的对 $(Y,j)$ 称为 $M$ 极化 K3 曲面。
若 K3 曲面的同构映射 $f\colon Y\to Y'$ 满足 $j=f^{*}\circ j'$，则称其为从 $(Y,j)$ 到 $(Y',j')$ 的同构映射。
当 $M$ 极化 K3 曲面满足
\begin{equation}
  j(C(M)^{+})\cap\Pic(Y)^{++}\neq\emptyset. \tag{4.15}
\end{equation}
时，称它为\textbf{丰沛}（ample）的。
\end{definition}
% [Tutor Note: 原图在 $\neq\emptyset$ 之后印有句点，而本句以「を満たすとき」未完待续，此句点疑为原文误排；照录。]

固定一个原始嵌入 $i_M\colon M\hookrightarrow L$，把它的正交补格称为 $T$。$M$ 极化 K3 曲面的周期区域
\begin{equation}
  \mathcal{D}=\{[\Omega]\in\mathbb{P}(T_{\mathbb{C}})\mid(\Omega,\Omega)=0,\ (\Omega,\overline{\Omega})>0\} \tag{4.16}
\end{equation}
是由 $T_{\mathbb{R}}$ 的已定向的正定平面构成的 Grassmann 流形 $O(2,19-t)/SO(2)\times O(19-t)$。
它由两个连通分支 $\mathcal{D}^{+}$ 与 $\mathcal{D}^{-}$ 组成，其中每一个都是 IV 型的有界 Hermite 对称区域。

令
\begin{equation}
  \Gamma(M)=\{\sigma\in O(L)\mid\text{对任意 }m\in M\text{ 有 }\sigma(m)=m\} \tag{4.17}
\end{equation}
并把 $\Gamma(M)$ 在自然单射 $\Gamma(M)\hookrightarrow O(T)$ 下的像记作 $\Gamma$。
由 K3 曲面的全局 Torelli 定理与周期映射的满射性，丰沛的 $M$ 极化 K3 曲面的模空间 $\mathcal{M}$
可以表示为 $\mathcal{D}^{\circ}/\Gamma$。这里，
\begin{equation}
  \mathcal{D}^{\circ}=\mathcal{D}\setminus\Bigl(\bigcup_{\delta\in\Delta(T)}H_{\delta}\cap\mathcal{D}\Bigr) \tag{4.18}
\end{equation}
是镜射面 $H_{\delta}=\{z\in T_{\mathbb{C}}\mid(z,\delta)=0\}$ 的补集。
把周期区域在其紧对偶
\begin{equation}
  \check{\mathcal{D}}=\{[\Omega]\in\mathbb{P}(T_{\mathbb{C}})\mid(\Omega,\Omega)=0\} \tag{4.19}
\end{equation}
中的闭包记作 $\mathcal{D}^{*}$。其拓扑边界为
\begin{equation}
  \mathcal{D}^{*}\setminus\mathcal{D}=\bigcup_{I\,:\,M_{\mathbb{R}}\text{ 的迷向子空间}}B(I) \tag{4.20}
\end{equation}
这里，$B(I)=\mathbb{P}(I_{\mathbb{C}})\setminus\bigl(\bigcup_{J\subsetneq I}\mathbb{P}(J_{\mathbb{C}})\bigr)$。
由于 $M$ 的符号为 $(2,19-t)$，$I$ 的秩为 1 或 2，而 $\mathbb{P}(I_{\mathbb{C}})\cap\mathcal{D}^{*}$ 或为一点，或为上半平面。
当 $I$ 在 $\mathbb{Q}$ 上有定义时，称 $B(I)$ 是有理的。
% [Tutor Note: 原书 p.132 的 (4.20) 下标印作「$I:M_{\mathbb{R}}$ の等方部分空間」，下一句又印作「$M$ の符号が $(2,19-t)$」。按上下文这两处的 $M$ 都应为 $T$：周期区域 (4.16) 取在 $\mathbb{P}(T_{\mathbb{C}})$ 中，而 $T=M^{\perp}\subset L$ 的符号正是 $(2,19-t)$（$M$ 的符号为 $(1,t)$，见 (4.12) 前的设定）。此处照原图录作 $M$。]
$\mathcal{M}$ 的佐竹–Baily–Borel 紧化为
\begin{equation}
  \mathcal{M}^{*}=\Bigl(\mathcal{D}\cup\bigcup_{I\,:\,\text{有理}}\mathbb{P}(I_{\mathbb{C}})\cap\mathcal{D}^{*}\Bigr)/\Gamma. \tag{4.21}
\end{equation}

% ===== 原书 p.133 =====

假定对某个格 $N$ 有 $T=U\perp N$，考虑与迷向子空间 $\mathbb{Z}e\subset T$ 相伴的尖点附近。
选取 $U$ 的基 $\{e,f\}$，使其满足 $(e,e)=(f,f)=0$ 以及 $(e,f)=1$，
并把 $e$ 在 $O(T)$ 中的稳定化子记作 $\Gamma_{e}$。对 $N$ 的元素 $v$，定义 $T$ 的等距映射 $\varphi_{e,v}$ 为
\begin{equation}
  \varphi_{e,v}(x)=x-\left(\frac{1}{2}(v,v)(e,x)+(v,x)\right)e+(e,x)v \tag{4.22}
\end{equation}
对任意 $w\in N$，容易看出
\begin{equation}
  \varphi_{e,v}\circ\varphi_{e,w}=\varphi_{e,v+w}, \tag{4.23}
\end{equation}
以及
\begin{equation}
  \varphi_{e,v}(e)=e,\qquad\varphi_{e,v}(f)=-\frac{1}{2}(v,v)e+f+v,\qquad
  \varphi_{e,v}(w)=-(v,w)e+w \tag{4.24}
\end{equation}
成立，并且
\begin{equation}
  \varphi_{e,\bullet}\colon N\hookrightarrow O(T). \tag{4.25}
\end{equation}
给出群的单射。
由于 $\Gamma_{e}$ 的任意元素 $\phi$ 都能借助满足 $\phi(f)\equiv f+v\bmod\mathbb{Z}e$ 的 $v\in N$
与适当的 $\psi\in O(N)$ 写成 $\psi\circ\varphi_{e,v}$，可知
\begin{equation}
  \Gamma_{e}=O(N)\ltimes N \tag{4.26}
\end{equation}
成立。
借助映射
\begin{equation}
  v\mapsto\Omega=-\frac{1}{2}(v,v)e+f+v \tag{4.27}
\end{equation}
可以把式 (4.16) 的周期区域与管状区域
\begin{equation}
  \{v=v_{1}+\sqrt{-1}v_{2}\in N_{\mathbb{C}}\mid(v_{2},v_{2})>0\} \tag{4.28}
\end{equation}
等同起来。在这一等同之下，$\varphi_{e,u}\in N\subset\Gamma_{e}$ 的作用变为平移 $v\mapsto v+u$。
因此，$\mathcal{D}/\Gamma$ 的尖点附近局部同构于
\begin{equation}
  (N_{\mathbb{C}}/N)/O(N)^{+}\cong N_{\mathbb{C}^{\times}}/O(N)^{+} \tag{4.29}
\end{equation}
这里，$O(N)^{+}$ 是 $O(N)$ 中保持连通分支 $\mathcal{D}^{+}$ 的指数 2 的子群。
若 $N$ 的扇 $\Sigma$ 在 $O(N)^{+}$ 作用下不变，则与 $\Sigma$ 相伴的环面簇 $X_{\Sigma}$ 具有 $O(N)^{+}$ 的自然作用。
把尖点附近替换为商 $X_{\Sigma}/O(N)^{+}$ 的原点附近，就得到 $\mathcal{D}/\Gamma$ 的环面部分紧化。

% CHUNK c03 DONE pages p006–p007
